\section{Introduction}

A coding agent can run for an hour and finish nothing. It keeps issuing tool calls, reading files, editing,
reverting, and re-running the same command while the task stands exactly where it stood before.
Practitioners recognise the pattern immediately --- the repeated \code{grep}, the edit that undoes the
previous edit, the twenty minutes of verification after the job was already done --- but the runtimes that
host them mostly do not, and the rules they do apply (``stop after three identical tool calls'') alarm on
productive work about as often as on stagnation.

This paper asks which observable signals let an online monitor tell that a coding agent has stopped making
task-relevant progress without knowing the correct solution. The monitor runs at every step, sees the task
statement and everything the agent has done and observed, and is denied the reference patch, the hidden
tests, future steps and the final grade. We measure each signal family against independently annotated
windows rather than against final success.

\paragraph{Why not simply measure file change.}
The naive operationalisation of progress is ``the workspace changed''. It is wrong in both directions. An
agent that reads fifteen files and finally locates the function that raises the reported error has made real
progress with an unchanged workspace; an agent that edits five files in a loop, undoing each change, has
changed the workspace a great deal and advanced nothing. Progress is therefore defined through three
channels an engineer would use judging a colleague's work in progress: \emph{epistemic} (something new and
task-relevant became known), \emph{implementation} (a plausible artefact changed durably) and
\emph{verification} (objective evidence about correctness improved). Stagnation is the joint absence of all
three while activity continues.

\paragraph{What we measured.}
We instrumented \pnTraj{} trajectories from \pcorpusName{} (\pnTasks{} tasks, \pcorpusScaffolds{} scaffolds,
\pcorpusModels{} models) and computed \pnWindows{} online feature vectors across six families: action
repetition, relevance-blind novelty, task-grounded evidence, verification deltas, workspace dynamics and
rolling semantic redundancy. The annotation codebook was written before any monitor and produced
\pannWindowsTotal{} labelled windows over \pannTrajTotal{} trajectories, each re-read by a second reader
under reduced context so that label stability could be measured. Monitors are trained on task-disjoint
folds and tested on held-out tasks, both as window classifiers and as sequential alarms under a false-stop
budget.

\paragraph{What we found.}
\begin{enumerate}[leftmargin=1.4em,itemsep=2pt,topsep=3pt]
\item \textbf{Rolling semantic redundancy is the strongest single family} (\paucBFoursemantic{} global
ROC-AUC, \pwithinBFoursemantic{} within a run). It beats repetition, verification and workspace churn, and
matches the best combination we could build.
\item \textbf{Task grounding is not the missing ingredient.} Weighting newly observed entities by whether
they mention the task reaches \paucCOneevidence{}, below plain exact repetition at \paucBTworepThree{} and
well below the semantic baseline ($\Delta=\ppairCOneevvsBFoursemDelta{}$, clustered interval excluding
zero). It adds value only in combination (\paucCThreeevidSem{}).
\item \textbf{Stagnation is an episode, not a window.} Consecutive same-label windows merge into
\pregEpisodes{} stall episodes with a median length of \pregMedianSteps{} steps. Scored at that unit, and
against productive stretches at matched positions within the same run, the monitor separates
\pregPmPct{}\% of episodes where a run-position proxy separates \pregProxySep{} of \pregProxyTot{}.
\item \textbf{The score drifts, so it cannot simply be thresholded.} Information flow rises with run
progress (positive trend in 93\% of runs, correlation $r=0.52$ with step index), which defeats any fixed
threshold. Anchoring the reference to a run's own opening, before a stall can begin, removes the drift and
yields a working alarm.
\item \textbf{The resulting alarm meets a practical target}: \palDet{}\% of stall episodes detected at a
\palFaWin{}\% per-window false-alarm rate, median latency zero steps, and \pnestedDet{}\% detection under a
nested protocol that selects every parameter on held-out runs.
\end{enumerate}

\paragraph{Contributions.}
(i) A falsifiable definition of stagnation with a published codebook, \pannWindowsTotal{} adjudicated
windows and a measured label-stability statistic. (ii) A six-family comparison of cheap online signals
under one protocol, finding that the intuitive signal --- task-grounded evidence --- loses to a task-blind
one, and identifying the representation drift that prevents either from being alarmed on directly. (iii) A
calibrated detector needing no labels, no reference trajectories and no per-task tuning, with the
annotations, frozen feature matrices and scripts that regenerate every figure and number here.

\paragraph{Scope.}
We study Terminal-Bench~2.0, a corpus of short, self-contained terminal tasks, so the conclusions are
strongest for runs of tens of steps and weakest for multi-hour repository work.
Section~\ref{sec:limits} states what the labels cannot support.
